\documentclass[11pt]{article}
\usepackage{series}
\usepackage{booktabs,tabularx,longtable,array}
\usepackage{xcolor}

\newcommand{\MTT}{\mathrm{MTT}}
\newcommand{\SM}{\mathrm{SM}}
\newcommand{\PEW}{P_{\mathrm{EW}}}
\newcommand{\Profile}{\mathcal P}
\newcommand{\Inputs}{\mathcal I}
\newcommand{\Outputs}{\mathcal O}
\newcommand{\Data}{\mathcal D}
\newcommand{\Redundancy}{\mathcal G}
\newcommand{\Ledger}{\mathfrak L}

\title{\textbf{Parameters, Source Provenance, and Structural Falsifiability in Modal Triplet Theory}\\
\large An Auditable Identifiability Ledger}
\author{Peter Nero}
\date{Version 2, July 2026}

\begin{document}
\maketitle

\begin{abstract}
A parameter count is meaningful only relative to a model, a physical tier, a
choice of observables, a quotient by redundancies, and a rule separating
construction data from held-out tests. This paper rebuilds the parameter and
falsifiability account of Modal Triplet Theory (MTT) around that principle.
We distinguish structural assumptions, selected discrete data, continuous
construction primitives, measured source coordinates, nuisance and metrology
inputs, derived outputs, and uncertainty data. We prove that raw symbol
counts are not invariant under reparameterization or gauge redundancy, that
replay of data used in construction is not a held-out prediction, that loss
of a global right inverse is not by itself a falsifier, and that a scalar
coherence-capacity diagnostic cannot eliminate the full source ledger.
At the current profile-standard Standard Model tier, the auditable ledger
contains one shared electroweak construction primitive and fifteen measured
common-scheme source coordinates. The transported couplings, charged Yukawa
rows, Higgs row, covariance workspace, and finite operators are outputs of
that declared construction and are not counted again. This is a coherent
embedding and reuse result, not a strict reduction of the Standard Model
parameter set. The stronger no-knob program remains two of nine upgrades
closed. We give valid falsification capsules for exact algebraic results,
numerical certificates, source provenance, and held-out empirical
predictions. The result is an honest parameter-identifiability framework:
MTT has substantial structural economy, but strict numerical economy must be
demonstrated rather than inferred from notation.
\end{abstract}

\section*{Version 2 revision note}

\begin{description}[style=nextline]
\item[Supersedes.]
Version 1, DOI
\url{https://doi.org/10.5281/zenodo.18255574}.
\item[Reason.]
Version 1 treated all apparent parameters as if they collapsed into one
invariant margin, asserted that MTT had no freely tunable structural
parameters, and listed the absence of a global measurable right inverse as a
theory-level falsifier. It also described broad physical programs as already
derived.
\item[Resolution.]
Version 2 replaces those claims by a tier-relative source ledger. It counts
independent inputs only after conventions and redundancies are declared,
separates profile replay from held-out prediction, and assigns every
falsifier to a named claim, domain, and evidence contract.
\item[Retained content.]
The useful aim survives: a theory with shared source objects can be more
constrained than a collection of sector-by-sector fits, and failures of
cross-sector compatibility can be powerful tests.
\item[Open boundary.]
MTT has not yet derived the full measured Standard Model profile from a
selected zero-primitive source. The strict program remains open, and this
paper does not convert profile inputs into predictions.
\end{description}

\tableofcontents

\section{Why parameter counting needs a ledger}

Statements such as ``the theory has one parameter'' or ``the theory has no
knobs'' sound precise, but they are incomplete. A count changes if one
introduces redundant coordinates, changes basis, quotients a gauge orbit,
replaces one vector by its components, or treats measured data as fixed
background rather than as input. A serious count must therefore answer five
questions:
\begin{enumerate}
\item What mathematical model is being evaluated?
\item Which physical or proof tier is being claimed?
\item Which quantities enter before the outputs are computed?
\item Which redundancies and conventions have been quotiented?
\item Which data were withheld from construction and used only for testing?
\end{enumerate}

This distinction is standard in identifiability theory
\cite{BellmanAstrom1970,RaueEtAl2009}. Structural identifiability asks whether
ideal observations distinguish parameter values in the declared model.
Practical identifiability asks whether finite, noisy data constrain them well
enough. Neither question can be answered from the number of symbols printed
in a formula.

The same discipline is especially important in MTT. Projection, bundle
selection, a finite operator, a normalized reserve vector, and an effective
action can reuse common structure. That reuse is valuable. It does not imply
that every scalar coefficient has been selected by the structure.

\subsection{Three different claims that are often confused}

We will keep the following claims separate.

\paragraph{Representation economy.}
Several sectors are represented by one carrier, one group action, or one
operator architecture. This can be exact even when numerical coefficients
remain external.

\paragraph{Construction economy.}
Once a declared source profile is supplied, many rows are generated without
additional sector-local fitting. This is stronger than representation
economy, but the source profile remains an input.

\paragraph{Predictive economy.}
A source fixed before target data are examined produces held-out observables
with a declared uncertainty budget. This is the tier required for a strict
parameter-reduction or no-knob claim.

Current MTT has exact examples of the first claim and a profile-standard
closure of the second. The third remains incomplete
\cite{MTTSuperset2026,MTTRoadmap2026}.

\section{The typed source ledger}

\begin{definition}[Realization and evidence ledger]
\label{def:ledger}
For a claimed result at tier \(T\), define
\[
\Ledger_T=
\bigl(
\mathcal S_T,\mathcal B_T,\mathcal A_T,\Profile_T,\mathcal N_T,
\Redundancy_T,\Outputs_T,\mathcal U_T,\Data_T
\bigr).
\]
The entries are:
\begin{description}[leftmargin=2.8cm,style=nextline]
\item[\(\mathcal S_T\): structural assumptions]
Hilbert spaces, bundles, domains, regularity, locality, positivity,
hyperbolicity, action principles, and theorem hypotheses.
\item[\(\mathcal B_T\): selected discrete data]
Ranks, branch labels, finite groups, representations, projectors, incidence
patterns, and other non-continuous choices.
\item[\(\mathcal A_T\): construction primitives]
Continuous quantities supplied by the construction before target observables
are evaluated.
\item[\(\Profile_T\): measured source profile]
Empirical coordinates used to initialize, calibrate, or transport the model.
\item[\(\mathcal N_T\): nuisance and metrology data]
Scheme, scale, detector response, calibration, covariance policy, and
external unit anchors.
\item[\(\Redundancy_T\): redundancies and conventions]
Gauge transformations, basis changes, field redefinitions, and equivalent
coordinate descriptions.
\item[\(\Outputs_T\): outputs]
Rows obtained after the preceding entries are fixed.
\item[\(\mathcal U_T\): uncertainty record]
Covariance, truncation error, interval bounds, numerical tolerance, and model
discrepancy.
\item[\(\Data_T\): data-use map]
For every datum, whether it is construction, validation, or held-out test
data.
\end{description}
\end{definition}

This ledger separates ontology from bookkeeping. A theorem hypothesis is not
automatically a fitted parameter. A discrete branch choice contributes no
continuous dimension but remains a selection obligation. A measured
coordinate is an empirical input even when it appears only once. An output
is not counted again merely because it is displayed in several bases.

\begin{definition}[Tier-relative continuous input count]
\label{def:count}
Let \(\Inputs_T=\mathcal A_T\times\Profile_T\times\mathcal N_T\).
When these spaces are smooth near the selected point and the declared
redundancy group \(\Redundancy_T\) acts regularly, define
\[
N_{\mathrm{cont}}(T)
=
\dim\bigl(\Inputs_T/\Redundancy_T\bigr).
\]
The dimensions of construction primitives, measured profile coordinates,
and nuisance data should also be reported separately. A single total is
insufficient when those categories have different scientific meanings.
\end{definition}

\begin{proposition}[Raw symbol counts are not invariant]
\label{prop:symbol-count}
The number of named scalar symbols in a presentation is not an invariant
parameter count. Under a local diffeomorphic reparameterization of the
quotient \(\Inputs_T/\Redundancy_T\), the intrinsic dimension in
\cref{def:count} is unchanged, while the printed symbol count may increase
or decrease.
\end{proposition}

\begin{proof}
Replace one coordinate \(a\) by \((u,v)\) subject to \(v-u^2=0\). The
presentation now has two symbols, but the constrained set is still
one-dimensional. Conversely, a vector with \(n\) components may be denoted by
one symbol without becoming one-dimensional. Local diffeomorphisms preserve
dimension, and quotienting a regular redundancy removes orbit directions.
Therefore only the quotient dimension, together with the category ledger,
has invariant content.
\end{proof}

\begin{remark}
Singular points require a stratified or algebraic dimension statement rather
than a naive Jacobian rank. This does not rescue a raw count; it makes the
required declaration more precise.
\end{remark}

\section{Identifiability, replay, and prediction}

Let
\[
F_T:\Inputs_T/\Redundancy_T\longrightarrow \Outputs_T
\]
be the declared forward map. Parameter economy and predictive power are
different properties of \(F_T\).

\begin{definition}[Local identifiability]
A selected input class \([\theta_\ast]\) is locally structurally identifiable
from an observable set \(J\) if there is a neighborhood \(U\) such that
\[
F_{T,J}([\theta])=F_{T,J}([\theta_\ast]),\qquad [\theta]\in U,
\]
implies \([\theta]=[\theta_\ast]\).
\end{definition}

Full-column Jacobian rank is a useful sufficient local test in a regular
finite-dimensional setting, but it is not a global uniqueness theorem.
Practical identifiability additionally depends on uncertainty and
conditioning.

\begin{definition}[Prediction certificate]
\label{def:prediction}
A held-out prediction certificate consists of:
\begin{enumerate}
\item a source and branch selection rule fixed before the test data enter;
\item disjoint construction and test sets
      \(\Data_{\mathrm{build}}\cap\Data_{\mathrm{test}}=\varnothing\);
\item a forward map with scheme, scale, units, and conventions;
\item an uncertainty budget independent of the realized test residual; and
\item an archived comparison rule fixed before evaluation.
\end{enumerate}
\end{definition}

\begin{proposition}[Replay is not held-out prediction]
\label{prop:replay}
Suppose a target datum \(y\) is used to choose an input, branch, normalization,
or acceptance rule, and the same \(y\) is then reproduced by the forward map.
That reproduction is a fit, calibration, or profile replay. It is not a
held-out prediction in the sense of \cref{def:prediction}.
\end{proposition}

\begin{proof}
The data-use sets are not disjoint. The construction map depends on \(y\), so
the displayed agreement does not test the map on information absent from its
construction. The conclusion follows directly from item 2 of
\cref{def:prediction}.
\end{proof}

\begin{example}
If fifteen measured common-scheme coordinates are transported to eight rows
and their covariance, the eight rows are genuine outputs of the transport.
They are not independent predictions of the fifteen inputs. The calculation
can still be exact, reproducible, and useful: it proves consistency and
scheme transport at the declared profile tier.
\end{example}

\section{What coherence capacity does and does not count}

The corrected coherence-capacity formalism begins with a sourced,
dimensionless reserve vector
\[
r(x)=\bigl(r_1(x),\ldots,r_m(x)\bigr),
\qquad
C(x)=\min_i r_i(x),
\]
and, when needed, a separately declared metric clearance
\cite{MTTCapacity2026}. Each row records a different admissibility
obligation, with its own provenance, units, normalization, and tolerance.

\begin{proposition}[Bottleneck compression is not source reconstruction]
\label{prop:min-not-injective}
For \(m\geq2\), the map
\[
(r_1,\ldots,r_m)\longmapsto \min_i r_i
\]
is not injective. Therefore the scalar bottleneck capacity cannot, without
additional structure, reconstruct the reserve vector or eliminate the
parameters on which its rows depend.
\end{proposition}

\begin{proof}
For every \(a>1\), the vectors \((1,a)\) have the same minimum \(1\).
Thus infinitely many distinct reserve vectors share one bottleneck value.
\end{proof}

This simple fact corrects a central overstatement of Version 1. Many
conditions may be summarized by their weakest reserve, but the summary does
not identify their sources. Nor does \(-\nabla C\) become a force, \(C\)
become a conserved charge, or a reparameterization \(f(C)\) become physically
equivalent merely because it preserves the zero set.

\section{The current MTT parameter ledger}

\subsection{Declared profile-standard closure}

Current MTT closes embedded renormalized-Standard-Model equivalence at a
declared profile standard. The accepted calculation has twelve of twelve
obligations closed, uses one shared electroweak construction primitive, and
uses fifteen measured common-scheme source coordinates
\cite{MTTSuperset2026,MTTRoadmap2026}. Standard SM quantization is imported
on this branch rather than derived from abstract projection.

\begin{table}[ht]
\centering
\small
\begin{tabularx}{\textwidth}{@{}p{3.2cm}p{1.9cm}X@{}}
\toprule
Ledger class & Count & Current status at the profile tier \\
\midrule
Structural assumptions & not a scalar count &
Selected branch, finite carrier, action and QFT conventions, regularity,
scheme, and theorem domains must be declared. \\
Selected discrete data & finite, tier-dependent &
Ranks, finite branch data, representation and projector choices. They add
zero continuous dimension but require provenance. \\
Construction primitives & \(1\) &
One shared physical electroweak primitive \(\PEW\) at the adopted standard. \\
Measured source profile & \(15\) &
Common-scheme SM source coordinates with a declared covariance policy. \\
Nuisance/metrology & reported separately &
Scale, scheme, unit anchors, likelihood or covariance choice, numerical
tolerances, and any external calibration. \\
Derived outputs & not recounted &
Transported couplings, nine charged Yukawa rows, the Higgs row, finite
operator entries, threshold rows, and covariance workspace after the source
is fixed. \\
Held-out observables & claim-specific &
Only observables excluded from construction can support predictive-economy
claims. \\
\bottomrule
\end{tabularx}
\caption{The current input/output ledger. Counts are meaningful only at the
declared profile tier.}
\label{tab:current-ledger}
\end{table}

Consequently, the transparent continuous construction ledger at this tier is
\[
N_{\mathrm{build}}^{\mathrm{profile}}=1+15=16
\]
before separately counted nuisance and metrology quantities. This equation
does not claim that the number 16 is directly comparable with every published
Standard Model parameter count. Such comparisons depend on the reference
scale, renormalization scheme, flavor basis, neutrino assumptions, and which
experimental inputs are treated as fixed. It says exactly what enters this
MTT execution.

\subsection{What has nevertheless been reduced}

The ledger does record real structural achievements.
\begin{itemize}
\item Native selected bundle automorphisms give the faithful
      \((SU(3)\times SU(2)\times U(1))/\mathbb Z_6\) group without adding
      continuous gauge-group knobs.
\item The selected finite carrier, anomaly table, one-Higgs projector, and
      finite gauge/ghost spectrum rows are exact at their declared tiers.
\item The same profile supports gauge, flavor, Higgs, threshold, and finite
      operator constructions without separately fitting each displayed
      output row.
\item The fifteen-to-eight precision transport has a positive-definite
      covariance workspace and explicit cross-covariance entries.
\end{itemize}

These are compatibility, reuse, and representation-economy results. They
constrain how inputs can be assembled. They do not yet replace the measured
profile by a selected physical source.

\subsection{The strict tier}

The stronger program asks for the source geometry and action to emit the
electroweak primitive and measured gauge, Yukawa, mixing, Higgs, threshold,
and precision values before those observables are used. The current strict
upgrade ledger is two of nine closed. The zero-primitive electroweak source
and true held-out Standard Model value packet remain open.

\begin{theorem}[Current-ledger statement]
\label{thm:current-ledger}
At the adopted profile standard, the declared one-plus-fifteen continuous
source ledger is sufficient to execute the accepted embedded-SM forward map.
Removing the measured profile leaves no current theorem that emits all of its
coordinates from the selected MTT geometry and action. Therefore:
\begin{enumerate}
\item profile-standard equivalence is closed at its stated tier;
\item a strict no-knob derivation is not closed;
\item outputs of the profile map are not additional independent inputs; and
\item those outputs are not held-out predictions of the source coordinates.
\end{enumerate}
\end{theorem}

\begin{proof}
Items 1 and 2 are the current A04 and A05 ledger statuses. The construction
uses one shared primitive and fifteen measured source coordinates. Once these
are fixed, the accepted forward map emits its transported and finite rows, so
they are outputs rather than additional inputs. By \cref{prop:replay}, rows
whose source values entered construction cannot be reclassified as held-out
predictions.
\end{proof}

\section{What can legitimately falsify a claim}

Falsification must be typed as carefully as the parameter count. A failed
theorem hypothesis, a failed numerical certificate, a failed physical model,
and a failed universal theory claim are not interchangeable.

\begin{definition}[Falsification capsule]
\label{def:falsifier}
A falsification capsule for a claim \(K\) contains:
\[
\mathfrak F(K)=
(K,\Omega,H,S,V,E,\mathcal R),
\]
where \(\Omega\) is the domain, \(H\) the hypotheses, \(S\) the selected
source, \(V\) the verifier, \(E\) the admissible error or uncertainty budget,
and \(\mathcal R\) the rule assigning a failure to the theorem, implementation,
realization, or broader theory.
\end{definition}

\begin{proposition}[Scope rule for certificate falsification]
\label{prop:scope-rule}
Let a claim state that verified hypotheses \(H\) imply a conclusion \(K\)
within error \(E\). If \(H\) and the source identity are independently
verified but \(V\) rejects \(K\) beyond \(E\), then that certified claim fails.
If \(H\) is not established for the tested system, the test shows
non-applicability, not falsification of the conditional theorem.
\end{proposition}

\begin{proof}
The first statement is modus tollens applied to the certified implication.
The second follows because a conditional statement makes no assertion outside
its hypotheses.
\end{proof}

\subsection{Valid falsifier classes}

\paragraph{Algebraic and descent falsifiers.}
An asserted representation, cocycle, anomaly cancellation, global-group
quotient, or commuting diagram can be checked exactly. Failure falsifies that
specific construction.

\paragraph{Analytic and numerical certificate falsifiers.}
A claimed spectral gap, coercivity margin, contraction inequality, interval
enclosure, tail bound, or convergence estimate can be replayed. A rejected
certificate falsifies the claimed numerical theorem at its declared inputs.

\paragraph{Source-provenance falsifiers.}
If a row advertised as source-selected actually depends on target data,
unreported calibration, a changed branch, or an untracked convention, its
source theorem fails even if the numerical value happens to agree.

\paragraph{Cross-sector consistency falsifiers.}
When one source is genuinely fixed independently and several sectors are
computed from it, incompatible outputs can falsify the shared-source
realization. This is stronger than noticing that two separate fits disagree.

\paragraph{Held-out empirical falsifiers.}
After \cref{def:prediction} is satisfied, a residual outside the predeclared
uncertainty and model-discrepancy budget falsifies the prediction capsule.
Whether it also falsifies a larger theory depends on the claim's stated
ownership.

\section{Why absence of a right inverse is not a universal falsifier}

Version 1 listed nonexistence of a global measurable right inverse across an
admissibility collapse as an exact theory-level falsifier. That criterion is
not valid in this form.

\begin{proposition}[No-section does not imply inconsistency]
\label{prop:no-section}
The absence of a global right inverse or section of a projection does not, by
itself, imply that the projection, its domain, or a physical model using it is
inconsistent.
\end{proposition}

\begin{proof}
The Hopf fibration \(S^3\to S^2\) is a well-defined smooth principal
\(U(1)\)-bundle with no global continuous section; a section would trivialize
the bundle, contradicting its nonzero first Chern class
\cite{BottTu1982}. Thus a mathematically consistent and physically useful
projection can lack a global section.
\end{proof}

The correct falsifier is conditional. If an MTT claim requires a section,
decoder, or inverse on a named domain, then a proof that no such object exists
falsifies that claim. If the model instead represents intentional information
loss, local trivializations, a quotient, or a restricted recovery channel,
the absence of a global inverse may be expected.

The same point applies to measurement, horizons, and projection-induced
irreversibility. Noninvertibility alone does not derive a physical arrow of
time, an outcome law, or a horizon. Those require a selected dynamics,
state, observable algebra, and operational domain.

\section{A current falsification matrix}

\begin{longtable}{@{}p{2.9cm}p{3.8cm}p{4.1cm}p{3.2cm}@{}}
\caption{Claim-level falsification contracts.}
\label{tab:falsifiers}\\
\toprule
Claim & Required source and domain & Valid failure test & Invalid overreach \\
\midrule
\endfirsthead
\toprule
Claim & Required source and domain & Valid failure test & Invalid overreach \\
\midrule
\endhead
Profile-standard embedded SM equivalence &
Frozen one-plus-fifteen ledger, branch, SMDR version, conventions, and
covariance policy &
Independent replay fails one of the twelve declared obligations or the
transport/covariance checks &
Calling a moving experimental central value a theorem target \\
\addlinespace
Native finite gauge structure &
Selected finite carrier, automorphism action, center quotient, and anomaly
table &
Exact matrix, representation, center, or anomaly calculation fails &
Inferring measured relative gauge couplings from normalized finite spectra \\
\addlinespace
Finite HYM or operator certificate &
Exact source hashes, discretization, norm, tail estimate, and tolerance &
Independent interval or a posteriori verifier rejects the claimed bound &
Treating one certified representative as global branch uniqueness \\
\addlinespace
Coherence-capacity certificate &
Complete reserve rows, scales, provenance, metric, and region &
A declared margin is nonpositive where uniform positivity was claimed &
Treating an arbitrary gradient of capacity as a force law \\
\addlinespace
Strict no-knob Standard Model values &
Source fixed before observed values, same-branch action, forward map, and
held-out data split &
Predicted held-out packet lies outside its predeclared uncertainty budget &
Counting successful profile replay as this strict test \\
\addlinespace
QM, QFT, gravity, or cosmology bridge &
Target-specific state, algebra, dynamics, observables, source map, and error
contract &
One required map or physical axiom fails on the claimed domain &
Falsifying all of MTT because an interpretive analogy is incomplete \\
\bottomrule
\end{longtable}

\section{Comparison with Standard Model and effective-theory practice}

The Standard Model is usually parameterized at a chosen scale and scheme,
with basis conventions and assumptions about neutrino masses. A numerical
comparison between MTT and the Standard Model is therefore meaningful only
after both ledgers use the same convention.

At present, MTT should make two separate statements.
\begin{enumerate}
\item \textbf{Structural statement.} One selected finite architecture
supports a faithful gauge group, chiral carrier, anomaly cancellation,
finite operator organization, one-Higgs projection, and exact finite
gauge/ghost spectra at their declared tiers.
\item \textbf{Numerical statement.} The accepted profile-standard execution
uses one shared construction primitive and fifteen measured source
coordinates. No strict global reduction of the measured Standard Model
profile is presently proved.
\end{enumerate}

This is not a negligible result. Structural sharing can expose consistency
conditions, remove duplicated fit slots, and provide a clear route toward
prediction. But credibility increases when the profile coordinates remain
visible in the ledger. Hiding them behind an operator or calling them
``geometry'' would not reduce their empirical role.

\section{Minimum standard for future parameter claims}

Every future claim of a derived constant, reduced parameter count, or
falsifiable prediction should publish the following record.

\begin{enumerate}
\item \textbf{Tier and claim owner.} Exact, numerical-certified, profile,
conditional, support, or open.
\item \textbf{Source identity.} Repository commit, artifact hash, branch, and
selected geometric or operator object.
\item \textbf{Input ledger.} Structural assumptions, discrete choices,
continuous primitives, measured coordinates, nuisance quantities, and unit
anchors.
\item \textbf{Quotient and convention map.} Gauge, basis, scheme, scale, and
field redefinitions.
\item \textbf{Forward map.} The executable map from inputs to outputs.
\item \textbf{Data-use certificate.} Construction, validation, and held-out
sets.
\item \textbf{Uncertainty budget.} Experimental covariance, truncation,
numerical, matching, and model-discrepancy errors.
\item \textbf{Independent verifier.} A replay path that does not trust only
the producer's success flag.
\item \textbf{Failure ownership.} What a failed test would reject and what it
would leave untouched.
\end{enumerate}

This standard prevents two opposite mistakes. It prevents an exact
structural theorem from being dismissed merely because a physical
normalization is still open. It also prevents a successful profile replay
from being promoted into a source theorem.

\section{What remains to be proved}

The strongest parameter-economy program now has a short, explicit agenda.

\begin{enumerate}
\item Derive \(\PEW\) from the selected source geometry and action
without an equivalent continuous physical primitive.
\item Emit the measured gauge, charged-Yukawa, mixing, Higgs, and threshold
source values from the same branch before those values enter construction.
\item Complete the nonperturbative QFT and renormalization/matching bridge
needed to compare physical observables with a controlled uncertainty budget.
\item Produce genuinely held-out predictions and archive the data-separation
certificate.
\item Compare the final MTT and Standard Model ledgers in one common scheme,
including discrete assumptions and nuisance data rather than only continuous
headline counts.
\end{enumerate}

These tasks are represented by the open strict-upgrade ledger and the
zero-primitive and true-precision blockers. Until they close, ``no-knob MTT''
is a research target, not a result.

\section{Version 2 changes and reasons}

\begin{table}[ht]
\centering
\small
\begin{tabularx}{\textwidth}{@{}p{3.2cm}p{4.6cm}X@{}}
\toprule
Version 1 statement & Version 2 decision & Reason \\
\midrule
MTT already derives QM, QFT, spacetime, irreversibility, and cosmology &
Withdraw as a global status statement &
Current bridges have different exact, profile, conditional, and open tiers. \\
All apparent parameters collapse to one invariant margin &
Withdraw and replace by the typed ledger &
The minimum of a reserve vector is noninjective and does not reconstruct its
sources. \\
No freely tunable structural parameters &
Replace by a tier-relative count &
The current profile execution declares one construction primitive and fifteen
measured source coordinates. \\
Persistence of an open degree of freedom falsifies MTT &
Narrow to failure of a declared source theorem or held-out prediction &
An open research obligation is not a failed theorem. \\
Absence of a global right inverse is a structural falsifier &
Withdraw as universal criterion &
Consistent nontrivial bundles and quotient maps may have no global section. \\
Many falsifiers follow automatically from few parameters &
Replace by falsification capsules &
Every test needs a domain, source, assumptions, verifier, error budget, and
failure owner. \\
\bottomrule
\end{tabularx}
\caption{Current-version delta.}
\end{table}

\section{Conclusion}

The honest current picture is sharper than either ``MTT has no parameters''
or ``MTT is merely another fit.'' MTT has exact structural results and a
profile-standard construction in which one architecture and one source
ledger feed many sectors. That reuse is a real achievement. The present
execution nevertheless uses one shared continuous construction primitive and
fifteen measured source coordinates, plus separately declared conventions,
metrology, and uncertainty data.

Falsifiability begins when a claim is made precise. Exact matrix and descent
claims can fail exactly. Numerical certificates can fail under independent
replay. Source claims can fail through hidden target dependence. Shared-source
models can fail through cross-sector incompatibility. Held-out predictions
can fail outside a predeclared uncertainty budget. The absence of a global
right inverse, the existence of an open parameter, or a diagnostic capacity
gradient is not by itself such a failure.

This ledger gives the research program a stable baseline. Future reductions
can now be measured as actual movements of quantities from the input columns
to the output columns, with no relabeling and no loss of provenance.

% BEGIN MTT MANAGED COMPUTATIONAL EVIDENCE
\section*{Computational Evidence and Reproducibility}

The current MTT status statements in this paper are imported from the curated
results repository. They are not rederived here. The source is frozen to:
\begin{quote}\small
\textbf{Repository:}
\url{https://github.com/PeterNero/mtt-results-repro}\\
\textbf{Commit:}
\texttt{31247ebb 5c22f3fb b5443024 365433c6 ee0bff4a}\\
\textbf{Manifest:} \path{release/result_manifest.json}\\
\textbf{Manifest SHA-256:}\\
\texttt{fb399689 60b00584 631dbf53 1a708e18}\\
\texttt{ef928d6b 6d935119 c185d7f6 32b1e7cd}
\end{quote}

\paragraph{Profile rows used directly.}
\begin{itemize}
\item \path{A01/current_global_lock}: current non-looping status and source
certificate map.
\item \path{A04/final_12_of_12_audit}: twelve-obligation profile-standard
embedded-SM audit.
\item \path{A02/precision_15_source_transport}: fifteen measured source
coordinates and transport workspace.
\item \path{A02/precision_8x8_workspace}: eight output coordinates and
positive-definite covariance.
\item \path{A01/charged_yukawa_higgs_profile}: charged Yukawa and Higgs
profile rows.
\end{itemize}

\paragraph{Exact structural rows used directly.}
\begin{itemize}
\item \path{A01/strict_pew_row}: exact source row at the declared
one-shared-primitive standard.
\item \path{A47/native_gauge_group}: native faithful gauge group and
parameter-assumption audit.
\item \path{A62/su3_finite_gauge_spectrum}: final exact finite gauge-spectrum
row and the common-normalized-spectrum no-go.
\end{itemize}

\paragraph{Open boundary.}
\begin{itemize}
\item \path{A05/strict_upgrade_ledger}: current \(2/9\) strict no-knob
upgrade ledger. This row is not evidence of strict closure.
\end{itemize}

Tier labels and status are inherited from the manifest. No imported row
promotes a profile value into a held-out prediction or changes theorem
ownership.
% END MTT MANAGED COMPUTATIONAL EVIDENCE

\bibliographystyle{plain}
\bibliography{main}

\end{document}
